% Consolidated guide-ch04 from 9a1384ddf1afcd855b09256c9c6f1c90b5023140; source fragments preserved.

% === BEGIN guide-ch04-examples.tex blob 4fc73bf20a1d0207bb834bf82dfde05d908c580e ===
% 《大学物理学学习指导》（第5版）第4章相对论：例4.1—4.5。
% 印刷页42—44 / 扫描PDF页50—52；按源图及文字独立推导。
% 非官方研究题解，待独立复核。
\sourceexample{4.1}{42--43}{50--51}
\problemtext 惯性系$S'$相对$S$沿$x$轴匀速运动。两事件在$S'$中发生于相同$x'$、不同$y'$处，且在$S'$中同时发生。问在$S$中是否仍同时。
\workings
标准洛伦兹变换取$S'$以速度$u$沿$S$的$+x$方向运动：
\[
\Delta t=\gamma\left(\Delta t'+\frac{u\Delta x'}{c^2}\right),
\qquad
\gamma=\frac1{\sqrt{1-u^2/c^2}}.
\]
题设给出$\Delta t'=0$、$\Delta x'=0$，故
\[
\boxed{\Delta t=0}.
\]
\crosscheck 虽然$\Delta y'\ne0$，横向坐标差不会出现在同一方向的洛伦兹时间变换中；这并不意味着沿$x'$有距离的同时事件在各系都同时。

\sourceexample{4.2}{43}{51}
\problemtext 一根静止于$S'$中、固有长度$l_0=1\,\mathrm m$的米尺，在$S'$中与$x'$轴成$30^\circ$，在$S$中测得与$x$轴成$45^\circ$。求$S'$相对$S$的速率和$S$中的尺长。
\workings
垂直于运动方向的尺寸不收缩：
$\Delta y=l_0\sin30^\circ=0.5\,\mathrm m$。
在$S$系中$\tan45^\circ=\Delta y/\Delta x$，故$\Delta x=0.5\,\mathrm m$。
原系中的纵向长度$\Delta x'=l_0\cos30^\circ=\sqrt3/2\,\mathrm m$，
由长度收缩
\[
\Delta x=\Delta x'\sqrt{1-\beta^2},\qquad \beta=\frac uc.
\]
所以
\[
\sqrt{1-\beta^2}=\frac{1}{\sqrt3},
\quad
\boxed{u=c\sqrt{\frac23}},
\qquad
\boxed{l=\sqrt{(\Delta x)^2+(\Delta y)^2}
=\frac1{\sqrt2}\,\mathrm m\approx0.7071\,\mathrm m}.
\]
\crosscheck 横向投影$0.5\,\mathrm m$不变，纵向投影从$0.8660$ m减为$0.5$ m，因此尺的倾角增大；与原书数值一致。

\sourceexample{4.3}{43--44}{51--52}
\problemtext 在$S$中，两事件发生于同一空间位置，时间间隔$\Delta t=2\,\mathrm s$；在沿$x$方向相对运动的$S'$中两事件时间间隔$\Delta t'=3\,\mathrm s$。求$S'$中两事件的空间距离。
\workings
因$\Delta x=0$，由洛伦兹变换
\[
\Delta t'=\gamma\Delta t,\quad
\gamma=\frac{3}{2}
\quad\Longrightarrow\quad
\beta=\sqrt{1-\frac1{\gamma^2}}=\frac{\sqrt5}{3}.
\]
位移变换$\Delta x'=\gamma(\Delta x-u\Delta t)$，因此
\[
\boxed{|\Delta x'|
=\gamma u\Delta t
=\frac32\cdot\frac{\sqrt5}{3}c\cdot2\,\mathrm s
=\sqrt5\,c\,\mathrm s}.
\]
取$c=3.00\times10^8\,\mathrm{m/s}$，得
$\boxed{|\Delta x'|\approx6.71\times10^8\,\mathrm m}$。
\crosscheck 不变量为
$c^2\Delta t^2-\Delta x^2
=c^2\Delta t'^2-\Delta x'^2$，即$4c^2=9c^2-5c^2$（时间按秒数代入）。相对运动符号取决于坐标轴正向，题问距离故取模。

\sourceexample{4.4}{44}{52}
\problemtext 一均匀直杆固有质量为$m_0$、固有长度$l_0$，沿杆长方向以匀速运动时在实验室系的测量长度为$l$。求整根杆的相对论动能。
\workings
长度收缩给$l/l_0=\sqrt{1-\beta^2}=\gamma^{-1}$。
相对论动能为总能量与静能之差，
\[
E_k=(\gamma-1)m_0c^2
\quad\Longrightarrow\quad
\boxed{E_k=m_0c^2\left(\frac{l_0}{l}-1\right)}.
\]
\crosscheck $l=l_0$时动能为零；$l<l_0$时动能为正。该式假设只计算刚体整体质心的平动动能，不包含因加速造成的内应力或温升。

\sourceexample{4.5}{44}{52}
\problemtext 两粒子静止质量均为$m_0$，在公共惯性系中沿同一直线以等大反向的速率$0.5c$运动。求每个粒子的动量大小与总能量，并在以其中一个粒子为静止参考系中求另一个粒子的动量和总能量。
\workings
在公共系中$\gamma(0.5)=1/\sqrt{1-0.5^2}=2/\sqrt3$，
\[
\boxed{|p|=\gamma m_0v=\frac{m_0c}{\sqrt3}\approx0.57735m_0c},
\qquad
\boxed{E=\gamma m_0c^2=\frac2{\sqrt3}m_0c^2\approx1.1547m_0c^2}.
\]
设换到第一个粒子的静止系，另一个粒子的速率由相对论速度合成得到
\[
v_{\rm rel}=\frac{0.5c+0.5c}{1+0.25}=0.8c,
\quad
\gamma(0.8)=\frac53.
\]
所以
\[
\boxed{|p'|=\frac53m_0(0.8c)=\frac43m_0c},
\qquad
\boxed{E'=\frac53m_0c^2}.
\]
\crosscheck 源书的约数分别为$0.58m_0c,1.15m_0c^2,1.33m_0c,1.67m_0c^2$。相对论动量必须使用$\gamma m_0v$，不能代入牛顿式$m_0v$。
% === END guide-ch04-examples.tex ===
